% Source chapter: market-rl (chapters/market-rl.tex, "Market-based architectures in RL and beyond"); moved verbatim from its Introduction.
\subsection{Related work}
\label{mba/sec:related}

\textbf{Market-based RL.} The majority of early work in this area has focused on market-based \emph{reinforcement learning} (RL) algorithms. The pioneering work in this domain consists of Holland's \emph{Learning Classifier Systems} or ``bucket brigade'' \citep{hollandPropertiesBucketBrigade1985a} in rule-based systems, where condition-action agents (``classifiers'') bid to post messages (which can be actions described in some language) onto a global message board. Improvements to this paradigm were made by Schmidhuber \citep{schmidhuberEvolutionaryPrinciplesSelfreferential1987, schmidhuberLocalLearningAlgorithm1989} who allowed agents to determine their own bids and imposed credit conservation, and by Baum \citep{baumModelIntelligenceEconomy1999} who further strictly enforced property rights, resulting in a much more familiar set-up called the ``Hayek Machine''. The Hayek machine, which can be applied to any Markov Decision Process (MDP), was subsequently extended to POMDPs by \citet{kweeMarketBasedReinforcementLearning2001} by adding external memory, and more recently \citep{changDecentralizedReinforcementLearning2020} modified the framework to use Vickrey auctions and prove that a Nash equilibrium of the market produces a globally optimal policy.\todo{is this much description necessary?}

\textbf{Bounded rationality and markets.} Other, more recent work in this area includes: \emph{Logical Induction} \citep{garrabrantLogicalInduction2020}, an algorithm that assigns probabilities to mathematical sentences based on their prices in a prediction market that (roughly speaking) pays off when a sentence is proven; and the \emph{Boundedly Rational Inductive Agent} (BRIA) \citep{oesterheldTheoryBoundedInductive2023} which solves finite decision problems by assigning it to the highest-bidding trader, similar to in market-based RL. The key insight of these works is that markets are useful for modeling \emph{boundedly rational agents}. The main results of each work -- the fact that the logical inductor cannot by dominated by any polynomial-time trader, and the ``boundedly rational inductive agent criterion'' in the latter paper -- are specific and precise formulations of the Efficient Market Hypothesis \citep{neymanAlgorithmicBayesianEpistemology2024}.\todo{remove last sentence?}

\textbf{Markets and neural networks.} A specific equivalence between classifier systems and neural networks has been studied in the \emph{Neural bucket brigade} \citep{davisMappingClassifierSystems1988, schmidhuberLocalLearningAlgorithm1989}, although this does not consider backpropagation. A suggestive analogy between markets and backpropagation is discussed in \citet{wentworthCompetitiveMarketsDistributed2018}, though only for the case of a strictly sequential, unit-width market like that in Def~\ref{mba:def:deep}. In our work we make this more precise and generalize it to \ref{mba:def:goods}\todo{maybe mention Agoric Computation, BitTensor and market-based scaffolding for LLMs? see comments}
%\textbf{Decentralized software} Also worth a mention are market-based ``decentralized intelligence'' protocols such as Agoric Computation \citep{millerMarketsComputationAgoric1988,millerComparativeEcologyComputational1988, drexlerIncentiveEngineeringComputational1988}, BitTensor \citep{raoBitTensorPeertoPeerIntelligence2021} and markets comprised of LLM agents \citep{liEconAgentLargeLanguage2024a}. AI economist paper? I thought it was about something different but people have told me it's related.
